% ============================================================
% DEBRIS MODEL JUSTIFICATION -- Team 0109D -- 04.4_Verification
% ============================================================
\documentclass[12pt, letterpaper]{article}
\usepackage[margin=0.5in]{geometry}
\usepackage[T1]{fontenc}
\usepackage[scaled=1]{helvet}
\renewcommand{\familydefault}{\sfdefault}
\usepackage[utf8]{inputenc}
\usepackage{booktabs}
\usepackage{amsmath}
\usepackage{enumitem}
\usepackage{hyperref}
\usepackage{xcolor}
\usepackage{titlesec}

\titleformat{\section}{\large\bfseries}{\thesection}{0.6em}{}
\titlespacing*{\section}{0pt}{8pt}{3pt}
\setlist{nosep, leftmargin=1.2em, itemsep=2pt}
\hypersetup{colorlinks=true, urlcolor=blue!60!black}
\setlength{\parskip}{4pt}
\setlength{\parindent}{0pt}

\begin{document}

\begin{center}
{\Large\textbf{Debris Analogue Selection Justification}}\\[2pt]
{\normalsize Team 0109D \quad$\vert$\quad ESC204 \quad$\vert$\quad Winter 2026}
\end{center}
\noindent\rule{\textwidth}{0.6pt}\vspace{6pt}

\section{Why We Needed a Debris Analogue}

We need to verify that the air blade can dislodge \textit{wet, adhered organic debris} from a surface that replicates asphalt shingles. Real autumn leaves at full scale don't behave the same way on a 35\,cm mock-up, so we needed a material that replicates the two things that actually make wet roof debris hard to clear:

\begin{enumerate}
    \item \textbf{Stiction when wet.} Wet leaves stick to shingles through surface tension and capillary forces at the contact interface. Whatever we use has to bond to a rough surface the same way when soaked.
    \item \textbf{Flat profile when saturated.} Wet leaves press flat against the roof. There are no protruding edges for air to get under. The analogue needs to go completely flat when wet too.
\end{enumerate}

\section{What We Used: Tissue Paper on 40-Grit Sandpaper}

\subsection*{The Surface: 40-Grit Aluminum Oxide Sandpaper}

40-grit sandpaper has a mean surface roughness around $R_a \approx 400$--$500\,\mu$m. Asphalt roofing shingles use embedded mineral granules (typically 0.5--1.5\,mm) that produce a similar rough, high-friction texture. The sandpaper gives us a comparable coefficient of static friction to what debris actually encounters on a real roof.

\subsection*{The Debris: Disaggregated Tissue Paper}

We separated tissue paper into individual thin layers and small pieces, then spread them across the 15\,cm-wide mock-up. We tested in two conditions:

\begin{itemize}
    \item \textbf{Dry:} Light fragments, easy to blow off. This just confirms the air source has enough raw force for loose debris, similar to dry leaves.
    \item \textbf{Saturated (wet):} Tissue soaked in water and pressed onto the 40-grit surface. Once wet, the tissue:
    \begin{itemize}
        \item Has \textbf{zero structural rigidity}: it conforms to every single grain of the sandpaper, maximizing contact area and capillary bonding
        \item Has \textbf{no aerodynamic profile}: lies completely flat with nothing sticking up for air to grab
        \item Creates \textbf{strong surface-tension adhesion}: the water film between tissue and sandpaper produces the same capillary stiction as rain between a wet leaf and a shingle
    \end{itemize}
\end{itemize}

\section{This Is Actually Harder Than Real Leaves}

Saturated tissue on 40-grit sandpaper is \textbf{harder to clear than real wet leaves on real shingles}:

\begin{table}[h]
\centering\small
\begin{tabular}{@{}lll@{}}
\toprule
\textbf{Property} & \textbf{Real Wet Leaf} & \textbf{Saturated Tissue on 40-Grit} \\
\midrule
Structural rigidity & Some (veins, stem) & None (conforms to every grain) \\
Contact area & Partial (leaf curvature) & Near-total (completely flat) \\
Aerodynamic profile & Low but nonzero & Effectively zero \\
Edges for air to grab & Some edges lift & No edges at all \\
\bottomrule
\end{tabular}
\end{table}

So if the air blade can peel saturated tissue off sandpaper at $d = 50$\,cm, it will definitely handle real wet leaves on real shingles. The tissue paper model is a deliberately conservative worst-case test.

\section{Results}

At $d = 50$\,cm on saturated tissue / 40-grit sandpaper:
\begin{itemize}
    \item \textbf{Circular Nozzle:} \textcolor{red}{\textbf{FAILED.}} The high-pressure column just pushed the flat, wet tissue harder against the surface. \href{https://drive.google.com/file/d/1wO-Wh8iFGygEvBvc3OOH91XQSwhoY_Ws/view}{\textit{[Video]}}
    \item \textbf{Custom Air Blade:} \textcolor{green!50!black}{\textbf{PASSED.}} The planar shear layer got underneath the water-sandpaper interface and peeled the debris off. \href{https://drive.google.com/file/d/16htpF3ptwKEM5UxYE679x3f3xWmt80h1/view}{\textit{[Video]}}
\end{itemize}

This confirms that the air blade's Coanda-effect shear mechanism can overcome worst-case stiction (Specifications O-1.1-STRU-1 and O-1.1-STRU-2).

\end{document}
